% ECONOMICS_PROBLEM: {"id": "C72-285", "title": "Statistical confidence in game equilibria", "jel_code": "C72", "source_id": "OP-0288"}
% !TeX program = xelatex
\documentclass[11pt,a4paper]{article}
\usepackage[margin=23mm]{geometry}
\usepackage{fontspec}
\setmainfont{Latin Modern Roman}
\usepackage{amsmath,amssymb,mathtools}
\usepackage{xcolor,enumitem,booktabs,longtable,array,xurl,hyperref}
\definecolor{muted}{HTML}{53636C}
\hypersetup{colorlinks=true,urlcolor=blue,linkcolor=blue}
\setlength{\parindent}{0pt}
\setlength{\parskip}{5pt}
\newcommand{\R}{\mathbb R}
\newcommand{\E}{\mathbb E}
\newcommand{\N}{\mathbb N}
\newcommand{\ind}{\mathbf 1}
\newcommand{\Var}{\operatorname{Var}}
\newcommand{\Cov}{\operatorname{Cov}}
\newcommand{\supp}{\operatorname{supp}}
\DeclareMathOperator*{\argmin}{arg\,min}
\DeclareMathOperator*{\argmax}{arg\,max}
\newcommand{\entrymeta}[3]{{\sffamily\small\color{muted}\textbf{\texttt{#1}}\quad|\quad #2\par #3\par}}
\newcommand{\Problemref}[1]{\texttt{#1}}
\newcommand{\JELref}[2]{\texttt{#2} (JEL \texttt{#1})}
\begin{document}
\section*{Statistical confidence in game equilibria}
\textbf{Problem ID:} \texttt{C72-285}\quad \textbf{JEL:} \texttt{C72}\par
% BEGIN ECONOMICS STATEMENT
\label{problem:OP-0288}
{\sffamily\small\color{muted}Original subject group: Game theory and strategic interaction\par}
\label{definitions:problem:30007570}
\textbf{Audit classification:} Published research agenda. The verified EGTA survey §§6 and 7.3 motivates statistical reasoning and exploration. Exact instance-optimal certificates are an editorial target; finite payoff estimation already gives a baseline.
\subsection*{Definitions and precise research target}
\textbf{Model and research target.} Let \(G\) be a finite \(n\)-player game with action sets \(A_i\) and unknown expected payoffs \(u_i(a)\in[0,1]\). A simulator queried at profile \(a\) returns independent bounded payoff samples with known query cost \(c(a)>0\); samples may have profile-dependent variance. An adaptive procedure chooses simulator queries, constructs candidate mixed profiles from the same observations, and stops at a random time with output \((\widehat\sigma,\widehat e)\). Define full-game exploitability
\[
 e_G(\sigma)=\max_i\max_{b_i\in A_i}
 \{u_i(b_i,\sigma_{-i})-u_i(\sigma)\}.
\]
For tolerances \(\varepsilon>0\) and \(\delta\in(0,1)\), require a simultaneously valid certificate \(\Pr(e_G(\widehat\sigma)\leq\widehat e\leq\varepsilon)\geq1-\delta\), uniformly over the maintained payoff class. Determine sharp expected simulation-cost bounds and adaptive sampling rules that approach them, accounting for candidate selection, optional stopping, and unevaluated deviations. State how bounds depend on action count, variance, payoff gaps, and exploitable structure. Exhaustively estimating every payoff is an available baseline; the research target is principled savings while preserving certification as strategy exploration changes the candidate. A valid solution must also specify failure or nontermination reporting when a prescribed computation budget is insufficient.

\emph{Definitions and maintained assumptions (editorial unless a source definition is named).} At profile \(a\), a query returns \(Y(a)\in[0,1]^n\) with mean \(u(a)\), independent of previous queries conditional on the adaptively chosen profile. Within-query player components may be correlated. Query cost is \(c(a)>0\). A stopping time \(\tau\), product mixed profile \(\widehat\sigma\) and bound \(\widehat e\) are measurable with respect to the data and procedure randomness. Demand coverage \(\Pr(e_G(\widehat\sigma)\le\widehat e)\ge1-\delta\) despite selection and stopping, and require success \(\widehat e\le\varepsilon\) almost surely at finite \(\tau\) for a guaranteed-completion procedure. A finite-budget procedure may instead output FAIL; coverage is assessed on all issued certificates, and its success probability is separately reported. Uniformly estimate every payoff to tolerance \(\eta\) and solve the estimated game to \(\xi\)-Nash: full exploitability is at most \(\xi+2\eta\) on the simultaneous estimation event. The unresolved target must impose an asymptotic size regime or informative structure and distinguish simulation cost from equilibrium computation cost.
\subsection*{Variants and assumption changes}
\begin{enumerate}
\item \textbf{Two-player zero-sum (editorial).} Use bounded sub-Gaussian payoff samples and known variances. Determine adaptive simulation-cost bounds for \(\varepsilon\)-saddle-point certification versus uniform sampling.
\item \textbf{General-sum exploration (editorial).} Jointly choose candidate strategies and deviation queries under optional stopping. Characterize when payoff gaps or low-dimensional structure yield savings while certifying all pure deviations; report failure when the budget cannot cover unseen actions.
\end{enumerate}
\subsection*{References and source audit}
\begin{enumerate}
\item §6 pp.1045–1052; §7.3 pp.1057–1058. Supports the stated research area; exact displayed specification is editorial. \url{https://arxiv.org/pdf/2403.04018}
\end{enumerate}
\begin{enumerate}\raggedright
\item\hypertarget{definitions:source:30007570:1}{} Michael P. Wellman; Karl Tuyls; Amy Greenwald. 2025. \emph{Empirical Game-Theoretic Analysis: A Survey}. §6 Statistical Reasoning in Empirical Games, pp.1045–1052; §7.3 Frontiers in Strategy Exploration, pp.1057–1058.
\url{https://arxiv.org/pdf/2403.04018}.
\end{enumerate}

\subsection*{Common conventions: Source mathematical conventions}

These conventions apply to each entry unless explicitly replaced.

\textbf{Models and identification.} A primitive model \(m\) specifies all probability laws, feasible choices, information, equilibrium and measurement rules used by its target. The admissible class \(\mathcal M\) is fixed independently of outcomes. If \(P_O(m)\) is its observable probability law and \(\tau(m)\) the target, its identified set at observable law \(P\) is
\[
 \mathcal I_\tau(P)=\{\tau(m):m\in\mathcal M,\ P_O(m)=P\}.
\]
Point identification means that this set is a singleton. An empty set signals incompatibility of data and assumptions. Coordinate bounds need not describe the joint set. Bounds are sharp only if every retained target is attainable or arbitrarily approachable by compatible models. Finite bounds require explicit restrictions on unobserved quantities; unrestricted confounding, hidden wealth, extrapolation or arbitrary equilibrium selection can yield uninformative sets.

\textbf{Causal interventions.} A potential outcome \(Y_i(p)\) is the outcome under a complete policy regime \(p\), including timing, financing, information and market spillovers. The target population and units are fixed before treatment. With interference, use the complete assignment vector or a justified exposure mapping; individual no-interference notation alone cannot identify an economy-wide intervention. A difference of mean potential outcomes does not require the cross-world joint distribution. Conditioning on post-treatment migration, survival, enrollment or employment changes the estimand and needs separate justification.

\textbf{Statistical statements.} An estimator, confidence set, forecast or test specifies a sampling law, model class and loss/coverage criterion. Uniform \((1-\alpha)\)-coverage means \(\inf_{m\in\mathcal M}\Pr_m\{\tau(m)\in C_n\}\geq1-\alpha\), or an explicitly asymptotic version. The whole parameter space trivially has coverage, so informativeness requires a stated width, risk or power objective. A forecasting improvement is relative to a named benchmark, information set and evaluation population.

\textbf{Equilibrium and optimization.} An equilibrium comprises feasible plans, optimality, beliefs where needed, budget constraints and all market/resource clearing. The primitives or a stated benchmark define each correspondence \(\mathcal E(p;m)\). Nonemptiness is assumed or proved, never inferred just from compact policy sets. A selected-equilibrium target uses a declared selection rule; a robust target quantifies over all permitted equilibria. Use suprema or infima unless attainment is established. An \(\varepsilon\)-optimal policy is feasible and within \(\varepsilon\) of the finite optimum value. Report an empty feasible set separately.

\textbf{Welfare and accounting.} Normative utility, interpersonal normalization, social weights, numeraire, discounting and the use of public receipts are primitives. Behavioral utility need not coincide with the welfare criterion. Household welfare includes ownership income and rents once; adding profits again to compensating variation generally double-counts them. A monetary equivalent is defined by an explicit transfer and price convention, with monotonicity and range conditions. Average effects, mechanism contributions, transfers and welfare losses are different targets. Infinite sums require convergence and dynamic feasible plans require terminal or no-Ponzi conditions.

\textbf{Source versions.} A working-paper date, revision date and journal publication year are distinguished. A DOI for a journal version is not automatically the DOI of an earlier manuscript. Locators refer to the stated version and printed pages where specified. A metadata-only check does not verify a claim about a full-text conclusion. Every entry's source subsection narrows the provenance claim accordingly.
% END ECONOMICS STATEMENT
\end{document}
